\addcontentsline{toc}{chapter}{Abstract}
Rapid Earth-observation response depends on two events occurring in the correct order: a spacecraft must receive a task, and it must then obtain a suitable access to the target before the deadline. This thesis studies whether assigning selected satellites a central-node role can improve that chain in a distributed satellite system. A central node is defined here as a preferred relay with broader link eligibility and an enhanced communication budget; it is not a constellation commander and it does not use a different sensor.

The experiment converts Sentinel-3 SLSTR active-fire detections into traceable wildfire follow-up tasks. It evaluates a full Walker-constellation grid that varies fleet size, plane count, altitude, inclination, and central-node fraction. Routing is replayed over finite, time-dependent communication windows, and an observation is credited only after the observing satellite has received the complete task. The post-processing separates mission performance, communication burden, failure response, and decision sensitivity.

The matched 60- and 120-satellite analysis shows that the denser India workload achieves 4.66 percentage points higher observation fulfilment than California on the same architectures, but lower complete useful-recipient dissemination and longer worst-case dissemination time. Increasing the central-node fraction improves some outcomes, especially relative to having no central-node links, but the response is not monotonic: high central-node fractions can reduce dissemination breadth under the bounded policy. Packet size is a major constraint; multiplying all packets by four reduces strict completion by 14.94 percentage points in California and 37.02 points in India. These results support a trade-space conclusion rather than a single universal optimum. A separate 180-satellite India sensitivity layer is interpreted conditionally because its task-provenance audit found a consistent but unintended 741-task population.

The main contribution is a reproducible link from wildfire data to communication-aware architecture evidence. The results show that orbital access alone is insufficient: routing participation, workload, contact capacity, and deadline timing jointly determine whether a potential observation becomes a completed mission outcome.
